\RevisionAddBegin
The requested application payload is the message length $q=|M|$. Gross mapping
capacity is $C_{\mathrm{gross}}$, auxiliary bits are the exact serialized wire
length $C_{\mathrm{aux}}=|A_{\mathrm{wire}}|_2$, and delivered net payload is
$P_{\mathrm{net}}=q-C_{\mathrm{aux}}$. The main text reports DRD
\cite{Lu2004DRD} as the primary distortion metric; changed-pixel and PSNR
outputs are retained in the reproducibility artifacts. Availability is the
fraction of test images that carry the target and complete the exact round trip.
\RevisionAddEnd

Primary per-payload comparisons use common available images. A secondary fixed
cohort contains the 46 images on which all four methods are available at all
four payloads. ABM is paired with each baseline for net bits and DRD using the
Wilcoxon signed-rank test and Holm correction within each payload. Mean net
payload is accompanied by a 95\% percentile bootstrap interval with 10,000
\RevisionAddBegin
image-level resamples and seed 20260608. Median DRD uses the analogous 95\%
percentile bootstrap interval computed from the same image-level resamples.
Runtime includes embedding,
\RevisionAddEnd
serialization, extraction, and verification. Memory is peak process RSS minus
pre-call RSS, sampled every 2 ms. Experiments used Python 3.13.5, NumPy 2.2.6,
PyTorch 2.8.0 CPU, scikit-learn 1.7.2, and SciPy 1.17.1 on 64-bit Windows with
12 logical processors.

\subsection{Serialization-Aware Activation Ablation}
\label{subsec:component-ablations}

Table~\ref{tab:activation-ablation} isolates the central algorithmic decision.
All variants use the same CNN distance-one candidate assignment; only active
table selection and exact wire charging differ. A0 reports gross payload before
wire cost and is therefore a gross-reporting diagnostic reference point. A1
charges the wire after activating all feasible tables. A2 is the proposed
serialization-aware activation, which selects the frequency-ordered prefix that
maximizes Equation~\ref{eq:net-capacity}. Relative to A1, A2 reduces auxiliary
information by 180 bits on average and raises mean net payload by 105.9 bits
(4.95\%) while preserving exact reversibility on all eight document images.

\begin{table*}[t]
\centering
\caption{Central ablation of serialization-aware mapping activation on the
eight document images. Values are means in bits except DRD; the post-hoc value
in A0 is the net payload obtained when the exact wire is charged after the
gross-only selection.}
\label{tab:activation-ablation}
\resizebox{\textwidth}{!}{
\begin{tabular}{@{}llcrrrrr@{}}
\toprule
\textbf{Variant} & \textbf{Active-table selection} & \textbf{Wire charged} &
\textbf{Gross cap.} & \textbf{Payload} & \textbf{Aux.} &
\textbf{Net payload} & \textbf{Mean DRD} \\
\midrule
A0 & all feasible; gross reported & Gross only & 3932 & 2948 & 0
& 2948 (2139 post hoc) & 0.771 \\
A1 & all feasible; wire charged & Charged & 3932 & 2948 & 809
& 2139 & 0.778 \\
A2 proposed & serialized-net prefix & Charged & 3832 & 2874 & 629
& 2245 & 0.761 \\
\bottomrule
\end{tabular}}
\end{table*}

\subsection{Candidate Ranking and CNN Architecture}

Table~\ref{tab:ranker-ablation} compares the learned ranker with simple
non-learning alternatives under the same admissible distance-one ZERO set and
common payload per image. The CNN is close to the direct DRD-cost ranker and
improves mean DRD by about 5\% relative to Hamming-index ordering. The
reversible mechanism remains the deterministic, globally injective PEAK/ZERO
substitution defined in
Section~\ref{sec:proposedmethod}.

\begin{table}[t]
\centering
\caption{Ranker ablation with identical distance-one candidates and common
payloads.}
\label{tab:ranker-ablation}
\begin{tabular}{@{}lrrr@{}}
\toprule
\textbf{Ranker} & \textbf{Mean net} & \textbf{Mean DRD} &
\textbf{Round trip} \\
\midrule
Hamming index & 2141 & 0.822 & Exact \\
Direct DRD cost & 2139 & 0.778 & Exact \\
Transition preserving & 2140 & 0.829 & Exact \\
Connectivity aware & 2140 & 0.820 & Exact \\
CNN & 2137 & 0.780 & Exact \\
\bottomrule
\end{tabular}
\end{table}

Table~\ref{tab:cnn-architecture-ablation} gives the controlled CNN
architecture diagnostic. All configurations use 20 epochs, AdamW with learning
rate $10^{-3}$, Smooth-L1 loss, batch size 128, the same two-image
train/validation split, and deterministic seeds. The deployed two-layer
w16 $9\times9$ model is retained as a compact compromise: it is close to the
lowest-MAE wider model, faster at inference, and keeps held-out placement DRD
within the same practical range.

\begin{table*}[t]
\centering
\caption{CNN architecture diagnostic on the two-image detailed split. Inference
time is measured on 4323 held-out candidates. Held-out DRD is the mean DRD
after model-ranked placement on the diagnostic held-out images. All
configurations share the reversible mapping and auxiliary stream.}
\label{tab:cnn-architecture-ablation}
\begin{tabular*}{\textwidth}{@{\extracolsep{\fill}}lrrrrrr@{}}
\toprule
\textbf{Configuration} & \textbf{Params} & \textbf{Train s} &
\textbf{Inf. $\mu$s} & \textbf{MAE} & \textbf{Pearson} &
\textbf{Held-out DRD} \\
\midrule
1 conv., $9\times9$, w24 & 14625 & 6.97 & 7.18 & 0.0716 & 0.962 & 0.848 \\
2 conv., $9\times9$, w24 & 38865 & 23.62 & 27.19 & 0.0339 & 0.991 & 0.830 \\
3 conv., $9\times9$, w24 & 44073 & 27.80 & 38.21 & 0.0347 & 0.991 & 0.848 \\
2 conv., $7\times7$, w24 & 38865 & 18.11 & 19.51 & 0.0383 & 0.989 & 0.845 \\
2 conv., $11\times11$, w24 & 38865 & 27.77 & 56.13 & 0.0485 & 0.985 & 0.825 \\
2 conv., $9\times9$, w16 (deployed) & 23649 & 15.16 & 24.86 & 0.0379 & 0.990 & 0.839 \\
2 conv., $9\times9$, w32 & 56385 & 27.27 & 40.62 & 0.0303 & 0.992 & 0.847 \\
\bottomrule
\end{tabular*}
\end{table*}

\subsection{Serialization Sensitivity for Baselines}

The original descriptions of PPOCP, Dong, and Huynh--Nguyen specify the
reversible principles while leaving the complete standalone byte stream partly
unspecified. The common protocol therefore uses conservative executable
synchronization. We also ran a header-credit sensitivity check: for each
baseline, the method-identifying header, version field, and image-shape fields
are removed from the charged cost, giving a 104-bit credit while retaining the
payload-critical maps, states, locations, and payload length. With this
credited accounting, ABM retains the highest mean net payload on every
all-method matched cohort: $-57$ vs.\ $-120$ bits at 64, $-7$ vs.\ $-72$ bits
at 128, $93$ vs.\ $23$ bits at 256, and $309$ vs.\ $200$ bits at 512 for the
best credited baseline. This fixed-header variation leaves the useful-capacity
ordering unchanged for the tested executable representations.

\subsection{Machine-Learning Validation}

Image-wise leave-one-out (LOSO) validation trains eight independent CNNs, each
on seven document images and tests it on the remaining complete image. Across
13,808 held-out candidates, pooled mean absolute error is 0.0367 and pooled
Pearson correlation is 0.9893; per-image correlations range from 0.9861 to
0.9939. A separate LOSO placement experiment applies each fold model to its
held-out image at an identical 75\% distance-one payload. Mean DRD decreases
from 0.8054 with deterministic Hamming-index ordering to 0.7663 with CNN
ordering, a 4.85\% reduction. All eight images improve, with per-image
reductions from 2.83\% to 7.94\%, while every message and cover is recovered
\RevisionAddBegin
exactly. The separate two-image architecture diagnostic in
Table~\ref{tab:cnn-architecture-ablation} evaluates 4323 held-out candidates
under the fixed 20-epoch protocol. Thus, regression accuracy and operational
DRD impact are evaluated separately, with fold-level regression and placement
results provided with the reproducibility artifacts.
\RevisionAddEnd

The uniform-block Random Forest is evaluated by source-image-grouped
cross-validation on 14,798 samples. It reaches ROC AUC 0.995, precision 0.897,
\RevisionAddBegin
recall 0.970, and flip-position accuracy 0.971. In the combined reversible
pipeline, the layer adds 181.4 gross bits per document image on average, while
its current coordinate representation uses 2691.0 auxiliary bits. The
resulting combined net payload is therefore $-264.6$ bits on average, compared
with 2245.0 bits for the base enumerative block-mapping layer. This identifies
coordinate coding as a concrete optimization opportunity and keeps the Random
Forest layer as a diagnostic component.
\RevisionAddEnd

\subsection{Document-Image and External Net Capacity}

Table~\ref{tab:document-net} reports the integrated CNN-ranked mapping layer at
a 75\% payload budget. Auxiliary sizes are measured from serialized bytes, and
all eight messages and covers are recovered exactly.

\begin{table}[t]
\centering
\caption{Net-aware results on the eight binary document images.}
\label{tab:document-net}
\begin{tabular}{@{}lrrrr@{}}
\toprule
\textbf{Image} & \textbf{Payload} & \textbf{Aux.} & \textbf{Net} & \textbf{DRD} \\
\midrule
circuit-2 & 1576 & 592 & 984 & 0.583 \\
formula-5 & 3692 & 672 & 3020 & 0.900 \\
graph1-5 & 1978 & 632 & 1346 & 0.662 \\
handwr2-8 & 2351 & 632 & 1719 & 0.768 \\
large-8 & 1138 & 488 & 650 & 0.616 \\
symbol-6 & 2367 & 672 & 1695 & 0.897 \\
table1-3 & 4161 & 656 & 3505 & 0.789 \\
french-4 & 5729 & 688 & 5041 & 0.864 \\
\midrule
\textbf{Average} & \textbf{2874} & \textbf{629} & \textbf{2245} & \textbf{0.760} \\
\bottomrule
\end{tabular}
\end{table}

\RevisionAddBegin
On the 100 BOSSbase test images, net-aware table selection gives a mean
embedded payload of 669.6 bits, 377.1 auxiliary bits, and 292.5 delivered net
bits. Median net capacity is 159 bits, 75 images are net positive, mean DRD is
0.300, and every evaluated message and cover is recovered exactly. Disabling
net-aware table selection on the same sample yields a mean embedded payload of
832.5 bits, 731.4 auxiliary bits, and 101.1 delivered net bits, showing that the
active-prefix decision also improves the external transfer setting.
\RevisionAddEnd

\RevisionAddBegin
\subsection{Frozen External Validation on FUNSD}
\label{subsec:funsd-external}

Table~\ref{tab:funsd-external} reports the frozen external validation on all
199 FUNSD scanned documents. The existing CNN ranker is used without retraining
or tuning. Two independently generated binary inputs are tested per source page:
fixed threshold 128 and Otsu. A1 activates all feasible tables and charges the
resulting wire, whereas A2 applies the proposed serialization-aware prefix
selection. Means are computed over available runs for each variant.
\RevisionAddEnd

\begin{table*}[t]
\RevisionAddBegin
\centering
\caption{Frozen external validation on 199 FUNSD scanned documents. Auxiliary
and net values are mean bits over available runs; DRD and positive-net counts
refer to A2. Availability is shown as available/199.}
\label{tab:funsd-external}
\resizebox{\textwidth}{!}{
\begin{tabular}{@{}lrrrrrrrrr@{}}
\toprule
\textbf{Binarization} & \textbf{$q$} & \textbf{A1 avail.} & \textbf{A2 avail.} &
\textbf{A1 aux.} & \textbf{A2 aux.} & \textbf{A1 net} & \textbf{A2 net} &
\textbf{A2 DRD} & \textbf{A2 net$>0$} \\
\midrule
Fixed 128 & 64  & 199/199 & 199/199 & 965.7 & 725.4 & $-901.7$ & $-661.4$ & 0.01049 & 0/199 \\
Fixed 128 & 128 & 199/199 & 199/199 & 965.7 & 725.4 & $-837.7$ & $-597.4$ & 0.02082 & 0/199 \\
Fixed 128 & 256 & 199/199 & 199/199 & 965.7 & 725.4 & $-709.7$ & $-469.4$ & 0.04124 & 1/199 \\
Fixed 128 & 512 & 198/199 & 198/199 & 969.3 & 727.8 & $-457.3$ & $-215.8$ & 0.08255 & 4/198 \\
\midrule
Otsu & 64  & 199/199 & 198/199 & 985.9 & 750.9 & $-921.9$ & $-686.9$ & 0.01015 & 0/198 \\
Otsu & 128 & 199/199 & 198/199 & 985.9 & 750.9 & $-857.9$ & $-622.9$ & 0.02008 & 0/198 \\
Otsu & 256 & 198/199 & 198/199 & 989.8 & 750.9 & $-733.8$ & $-494.9$ & 0.03981 & 0/198 \\
Otsu & 512 & 197/199 & 197/199 & 993.2 & 753.4 & $-481.2$ & $-241.4$ & 0.07928 & 3/197 \\
\bottomrule
\end{tabular}}
\RevisionAddEnd
\end{table*}

\RevisionAddBegin
Every available FUNSD run recovers the payload exactly and restores the binary
cover bit-for-bit. On paired available images, A2 is never worse than A1 in
net payload: the fraction $P_{\mathrm{net}}^{A2}\geq P_{\mathrm{net}}^{A1}$ is
1.0 at every payload under both binarizations. For fixed-128, mean A2 net gains
(and equal auxiliary savings) are 240.3, 240.3, 240.3, and 241.4 bits at 64,
128, 256, and 512 bits, respectively; for Otsu they are 238.9, 238.9, 238.9,
and 239.8 bits. The median gain is 232 bits in every condition. The associated
mean DRD differences $\mathrm{DRD}_{A2}-\mathrm{DRD}_{A1}$ remain small,
ranging from $6.66\times10^{-5}$ to $7.99\times10^{-4}$ across the eight
binarization/payload conditions.

The external result therefore supports two transfer claims that are distinct
from net-efficiency claims. First, exact reversible operation generalizes to a
large real scanned-document corpus without retraining. Second, the proposed
serialization-aware activation continues to reduce synchronization overhead on
every paired operating point. By contrast, the delivered net payload remains
metadata-dominated throughout the tested FUNSD range: mean net payload is
negative through 512 bits, and the largest observed positive-net fraction is
4/198 (2.02\%). Thus FUNSD validates reversibility and the A2 activation
benefit, but it does not establish positive net efficiency for scanned
documents at these payload sizes.
\RevisionAddEnd

\subsection{Multi-Payload Common-Protocol Results}

\RevisionAddBegin
Table~\ref{tab:multiload} and Figure~\ref{fig:multiload} summarize the
all-method matched BOSSbase subsets. ABM moves from metadata-dominated operation
at 64 bits to a near break-even 128-bit point: the payload-specific matched
cohort has a $-7$-bit mean, while the fixed 46-image cohort has an 8-bit mean.
It reaches positive mean net payload at 256 and 512 bits on these BOSSbase
cohorts. PPOCP defines the lowest-DRD curve, while ABM defines the highest-net
curve.
\RevisionAddEnd

